\begin{scanpage}{0139}{121}
Summing the secondary wave which has been found with the unperturbed primary wave from edge 1; we obtain
\[
\begin{aligned}
H_z(1-2)={}&\\
={}&H_{0z}\frac2\pi\widetilde g(1)\times\int_\infty^{2\sqrt{ka}\cos(\pi/4-\varphi/2)}e^{iq^2}\,dq\,\frac{e^{ikr}}{\sqrt{2kr}}e^{ika(\sin\alpha-\sin\varphi)}\\
&+H_{0z}g(1)\frac{e^{i(kr+\pi/4)}}{\sqrt{2\pi kr}}e^{ika(\sin\alpha-\sin\varphi)}.
\end{aligned}\tag{21.11}
\]
This expression reduces to zero if one assumes $\phi=-\pi/2$; consequently, the secondary diffraction eliminates the field discontinuities which occurred in the previous approximation when $\phi=-\pi/2$. However, in the direction $\phi=\pi/2$ the field (21.11) is different from zero. Since $H_z$ is an odd function of the $x$ coordinate, the relationship $H_z|_{\varphi=\pi/2-0}\ne0$ means that the fringing field components $H_z$ and $E_\phi$ will undergo a discontinuity with a transition through the direction $\phi=\pi/2$. The reason for such a jump, as before, is that in our calculations the plane $x=0$ is a plane of currents. By finding the secondary wave from edge 2, we actually considered that the diffraction takes place not on the edge of a finite width strip, but on the edge of an ideally conducting half-plane $-a<y<\infty$.

Again the resulting discontinuity has an order of magnitude of $\frac1{\sqrt{ka}\sqrt{kr}}$. It is clear that one may completely eliminate the field discontinuities only with consideration of multiple diffraction. However, the calculation of fields arising with multiple diffraction requires specific consideration of the following terms in order of smallness in the expansion of the primary edge in inverse powers of $\sqrt{kr}$ (see, for example, [46]). All this greatly complicates the calculations. Therefore, we, using the condition $ka\gg1$, will limit ourselves to an investigation of secondary diffraction, and in order to eliminate the discontinuities in the plane $x=0$, we will proceed in the following way.
\end{scanpage}
